\documentclass[10pt,a4paper]{amsart}
\usepackage{amsmath,amssymb,amsthm,graphicx}
\usepackage{verbatim}
\newenvironment{reference}{\comment}{\endcomment}
%\newenvironment{reference}{}{}
\newenvironment{slogan}{\comment}{\endcomment}
\newenvironment{history}{\comment}{\endcomment}

% For commutative diagrams we use Xy-pic
\usepackage[all]{xy}

% We use 2cell for 2-commutative diagrams.
\xyoption{2cell}
% Theorem environments.
%
\theoremstyle{plain}
\newtheorem{theorem}[subsection]{Theorem}
\newtheorem{proposition}[subsection]{Proposition}
\newtheorem{lemma}[subsection]{Lemma}

\theoremstyle{definition}
\newtheorem{definition}[subsection]{Definition}
\newtheorem{example}[subsection]{Example}
\newtheorem{exercise}[subsection]{Exercise}
\newtheorem{situation}[subsection]{Situation}

\theoremstyle{remark}
\newtheorem{remark}[subsection]{Remark}
\newtheorem{remarks}[subsection]{Remarks}

\numberwithin{equation}{subsection}

% Macros
%
\def\lim{\mathop{\mathrm{lim}}\nolimits}
\def\colim{\mathop{\mathrm{colim}}\nolimits}
\def\Spec{\mathop{\mathrm{Spec}}}
\def\Hom{\mathop{\mathrm{Hom}}\nolimits}
\def\Ext{\mathop{\mathrm{Ext}}\nolimits}
\def\SheafHom{\mathop{\mathcal{H}\!\mathit{om}}\nolimits}
\def\SheafExt{\mathop{\mathcal{E}\!\mathit{xt}}\nolimits}
\def\Sch{\mathit{Sch}}
\def\Mor{\mathop{\mathrm{Mor}}\nolimits}
\def\Ob{\mathop{\mathrm{Ob}}\nolimits}
\def\Sh{\mathop{\mathit{Sh}}\nolimits}
\def\NL{\mathop{N\!L}\nolimits}
\def\CH{\mathop{\mathrm{CH}}\nolimits}
\def\proetale{{pro\text{-}\acute{e}tale}}
\def\etale{{\acute{e}tale}}
\def\QCoh{\mathit{QCoh}}
\def\Ker{\mathop{\mathrm{Ker}}}
\def\Im{\mathop{\mathrm{Im}}}
\def\Coker{\mathop{\mathrm{Coker}}}
\def\Coim{\mathop{\mathrm{Coim}}}

% Boxtimes
%
\DeclareMathSymbol{\boxtimes}{\mathbin}{AMSa}{"02}

%
% Macros for moduli stacks/spaces
%
\def\QCohstack{\mathcal{QC}\!\mathit{oh}}
\def\Cohstack{\mathcal{C}\!\mathit{oh}}
\def\Spacesstack{\mathcal{S}\!\mathit{paces}}
\def\Quotfunctor{\mathrm{Quot}}
\def\Hilbfunctor{\mathrm{Hilb}}
\def\Curvesstack{\mathcal{C}\!\mathit{urves}}
\def\Polarizedstack{\mathcal{P}\!\mathit{olarized}}
\def\Complexesstack{\mathcal{C}\!\mathit{omplexes}}
% \Pic is the operator that assigns to X its picard group, usage \Pic(X)
% \Picardstack_{X/B} denotes the Picard stack of X over B
% \Picardfunctor_{X/B} denotes the Picard functor of X over B
\def\Pic{\mathop{\mathrm{Pic}}\nolimits}
\def\Picardstack{\mathcal{P}\!\mathit{ic}}
\def\Picardfunctor{\mathrm{Pic}}
\def\Deformationcategory{\mathcal{D}\!\mathit{ef}}
\usepackage{cleveref}
\crefname{originalreference}{original reference}{original references}
\Crefname{originalreference}{Original reference}{Original references}
\makeatletter
\newlabel{properties-lemma-quasi-coherent-finite-type-ideals}{{01PH}{source}{}{sourcepointer-1}{}}
\newlabel{properties-lemma-quasi-coherent-finite-type-ideals@cref}{{[originalreference][0][]01PH}{[1][0][]source}{}{}{}}
\makeatother
\crefname{originallemma}{original lemma}{original lemmas}
\Crefname{originallemma}{Original lemma}{Original lemmas}
\makeatletter
\newlabel{lemma-categorically-compact-QCoh}{{0GPE}{source}{}{sourcepointer-2}{}}
\newlabel{lemma-categorically-compact-QCoh@cref}{{[originallemma][0][]0GPE}{[1][0][]source}{}{}{}}
\makeatother
\crefname{originallemma}{original lemma}{original lemmas}
\Crefname{originallemma}{Original lemma}{Original lemmas}
\makeatletter
\newlabel{lemma-supported-on-support}{{0GPG}{source}{}{sourcepointer-3}{}}
\newlabel{lemma-supported-on-support@cref}{{[originallemma][0][]0GPG}{[1][0][]source}{}{}{}}
\makeatother
\crefname{originallemma}{original lemma}{original lemmas}
\Crefname{originallemma}{Original lemma}{Original lemmas}
\makeatletter
\newlabel{lemma-quotient-supported-on-closed}{{0GPH}{source}{}{sourcepointer-4}{}}
\newlabel{lemma-quotient-supported-on-closed@cref}{{[originallemma][0][]0GPH}{[1][0][]source}{}{}{}}
\makeatother
\crefname{originallemma}{original lemma}{original lemmas}
\Crefname{originallemma}{Original lemma}{Original lemmas}
\makeatletter
\newlabel{lemma-characterize-affine}{{0GPI}{source}{}{sourcepointer-5}{}}
\newlabel{lemma-characterize-affine@cref}{{[originallemma][0][]0GPI}{[1][0][]source}{}{}{}}
\makeatother
\crefname{originallemma}{original lemma}{original lemmas}
\Crefname{originallemma}{Original lemma}{Original lemmas}
\makeatletter
\newlabel{lemma-functor-equivalence}{{0GP0}{source}{}{sourcepointer-6}{}}
\newlabel{lemma-functor-equivalence@cref}{{[originallemma][0][]0GP0}{[1][0][]source}{}{}{}}
\makeatother
\crefname{originalreference}{original reference}{original references}
\Crefname{originalreference}{Original reference}{Original references}
\makeatletter
\newlabel{Brandenburg}{{shared-bibliography-key}{source}{}{sourcepointer-7}{}}
\newlabel{Brandenburg@cref}{{[originalreference][0][]shared-bibliography-key}{[1][0][]source}{}{}{}}
\makeatother
\title{Gabriel-Rosenberg reconstruction: equivalences of QCoh come from isomorphisms of schemes and invertible modules}
\author{The Stacks Project Authors}
\date{stacks-project:functors.tex:local-snapshot-20261012; GNU FDL 1.2 or later}
\begin{document}
\maketitle
\noindent\textit{Editorial note.} All delivered mathematical text is the author's own, copied verbatim from the pinned Stacks Project chapter file functors.tex; no statement or proof step was rewritten, reordered or removed.\par\medskip
\noindent\textit{Editorial note.} Span recomputation by label pairing (LANE\_BRIEF 10.2): in this chapter the author's proof environment follows the statement environment as a standalone block; the spec covers the complete environment pair (statement span, proof span) as recorded in the pool candidate.\par\medskip
\noindent\textit{Editorial note.} Occupation scan: combined claim\_key + statement\_anchor + generic-line scan over all registration files (active and withdrawn) found 0 hits for functors.tex (fu-occupied.json, 2026-10-14); consistent with the duty manager's six-check dispatch.\par\medskip
\ifdefined\hypertarget\hypertarget{sourcepointer-1}{}\fi
\noindent\textit{Source reference.} Original reference 01PH: sources/stacks-chow/tags.txt, source lines 2230--2230. Cross-chapter standard prerequisite cited by the author inside the counted span; it is used as a pointer and its own proof is not part of this item.\par\medskip
\ifdefined\hypertarget\hypertarget{sourcepointer-2}{}\fi
\noindent\textit{Source reference.} Original lemma 0GPE: sources/stacks-chow/tags.txt, source lines 20824--20824. Cross-chapter standard prerequisite cited by the author inside the counted span; it is used as a pointer and its own proof is not part of this item.\par\medskip
\ifdefined\hypertarget\hypertarget{sourcepointer-3}{}\fi
\noindent\textit{Source reference.} Original lemma 0GPG: sources/stacks-chow/tags.txt, source lines 20826--20826. Cross-chapter standard prerequisite cited by the author inside the counted span; it is used as a pointer and its own proof is not part of this item.\par\medskip
\ifdefined\hypertarget\hypertarget{sourcepointer-4}{}\fi
\noindent\textit{Source reference.} Original lemma 0GPH: sources/stacks-chow/tags.txt, source lines 20827--20827. Cross-chapter standard prerequisite cited by the author inside the counted span; it is used as a pointer and its own proof is not part of this item.\par\medskip
\ifdefined\hypertarget\hypertarget{sourcepointer-5}{}\fi
\noindent\textit{Source reference.} Original lemma 0GPI: sources/stacks-chow/tags.txt, source lines 20828--20828. Cross-chapter standard prerequisite cited by the author inside the counted span; it is used as a pointer and its own proof is not part of this item.\par\medskip
\ifdefined\hypertarget\hypertarget{sourcepointer-6}{}\fi
\noindent\textit{Source reference.} Original lemma 0GP0: sources/stacks-chow/tags.txt, source lines 20810--20810. Cross-chapter standard prerequisite cited by the author inside the counted span; it is used as a pointer and its own proof is not part of this item.\par\medskip
\ifdefined\hypertarget\hypertarget{sourcepointer-7}{}\fi
\noindent\textit{Source reference.} Original reference shared-bibliography-key: sources/stacks-functors/SOURCE.json, source lines 14--15. Author's bibliographic citation inside the counted span (reference environment pointing at Brandenburg); the shared Stacks bibliography is not staged locally, so the cite key is declared here per SOURCE.json bibliography\_note. Pointer only; its content is not part of this item.\par\medskip
\begin{proposition}
\label{proposition-gabriel-rosenberg}
\begin{reference}
Special case of \cite[Theorem 1.2]{Brandenburg}
\end{reference}
Let $X$ and $Y$ be quasi-compact and quasi-separated schemes.
If $F : \QCoh(\mathcal{O}_X) \to \QCoh(\mathcal{O}_Y)$
is an equivalence, then there exists an isomorphism
$f : Y \to X$ of schemes and an invertible $\mathcal{O}_Y$-module
$\mathcal{L}$ such that $F(\mathcal{F}) = f^*\mathcal{F} \otimes \mathcal{L}$.
\end{proposition}

\begin{proof}
Of course $F$ is additive, exact, commutes with all limits,
commutes with all colimits, commutes with direct sums, etc.
Let $U \subset X$ be an affine open subscheme.
Let $\mathcal{I} \subset \mathcal{O}_X$ be a finite type
quasi-coherent sheaf of ideals such that $Z = V(\mathcal{I})$
is the complement of $U$ in $X$, see
Properties, Lemma \ref{properties-lemma-quasi-coherent-finite-type-ideals}.
Then $\mathcal{O}_X/\mathcal{I}$ is a finitely presented
$\mathcal{O}_X$-module. Hence $\mathcal{G} = F(\mathcal{O}_X/\mathcal{I})$
is a finitely presented $\mathcal{O}_Y$-module by
Lemma \ref{lemma-categorically-compact-QCoh}.
Denote $T \subset Y$ the support of $\mathcal{G}$ and set
$V = Y \setminus T$. Since $\mathcal{G}$ is of finite presentation,
the scheme $V$ is a quasi-compact open of $Y$.
By Lemma \ref{lemma-supported-on-support} we see that $F$ induces an
equivalence between
\begin{enumerate}
\item the full subcategory of $\QCoh(\mathcal{O}_X)$ consisting
of modules supported on $Z$, and
\item the full subcategory of $\QCoh(\mathcal{O}_Y)$ consisting
of modules supported on $T$.
\end{enumerate}
By Lemma \ref{lemma-quotient-supported-on-closed} we obtain a commutative
diagram
$$
\xymatrix{
\QCoh(\mathcal{O}_X) \ar[r]_F \ar[d] &
\QCoh(\mathcal{O}_Y) \ar[d] \\
\QCoh(\mathcal{O}_U) \ar[r]^{F_U} &
\QCoh(\mathcal{O}_V)
}
$$
where the vertical arrows are the restruction functors and the
horizontal arrows are equivalences. By Lemma \ref{lemma-characterize-affine}
we conclude that $V$ is affine. For the affine case we have
Lemma \ref{lemma-functor-equivalence}. Thus we find that there
is an isomorphism $f_U : V \to U$ and an invertible
$\mathcal{O}_V$-module $\mathcal{L}_U$ such that
$F_U$ is the functor
$\mathcal{F} \mapsto f_U^*\mathcal{F} \otimes \mathcal{L}_U$.

\medskip\noindent
The proof can be finished by noticing that the diagrams above
satisfy an obvious compatibility with regards to inclusions
of affine open subschemes of $X$. Thus the morphisms $f_U$ and the
invertible modules $\mathcal{L}_U$ glue. We omit the details.
\end{proof}

\end{document}
